%% ============================================================
%% VALIDATION RESULTS -- v2, REWRITTEN 2026-08-16 after the model-attribution correction.
%% Drop-in for <<KEEP:validation-results>>.
%%
%% SUPERSEDES the earlier version of this file, which framed the temperature metrics as
%% "gates v1 reported as met and we now report as failed". That framing was wrong and has
%% been removed. v1's targets table assigned the supply-temperature gates to the
%% NLP formulation explicitly ("only this variant propagates supply temperature along
%% pipes"), and its results table scored them there. The attribution was correct in v1;
%% it was lost in the interim v2 table, which scored MILP output against NLP gates.
%% With the Model column restored, nothing in v1 needs to be walked back.
%%
%% What DOES remain different from v1, and is stated below: the NLP temperature figure in
%% v1 was a mean over six of fourteen nodes across a 744 h winter window; the figure
%% reported here is the far-end node over the full annual record, unselected. Both are
%% true. The appendix per-node table shows the selection was not favourable -- the all-node
%% mean sits between the group means -- so this is a change of view, not a correction.
%% ============================================================

Table~\ref{tab:validation} reports the validation outcome against the gates fixed ex ante,
with each metric scored against the formulation it applies to. That attribution is
essential to reading the table: the baseline prescribes supply temperature from the heating
curve and does not propagate it along pipes, so downstream temperature metrics are
structurally outside its scope and are marked as such rather than as failures.

\paragraph{What the baseline is validated on.}
Three quantities test the formulation that produces every cost result in this paper, and
all three pass. The annual delivered energy is reproduced to \textbf{1.23}\,\% against a
$2$\,\% gate. This is the decisive one: the network loss is the difference between generated
and delivered heat, so an annual energy match at this level fixes the annual loss at the
same level, and the loss is what the entire decomposition prices. The return temperature at
the source matches to \textbf{2.08}\,°C against a $2.5$\,°C gate, and -- more informative
than the magnitude -- with a bias of only $-0.31$\,°C. The source return sensor is the one
temperature instrument in the network that does \emph{not} sit downstream of a mixing valve,
and against it the model is essentially unbiased. Heat-pump performance sits in its expected
band at a mean coefficient of performance of \textbf{2.99}.

\paragraph{What cannot be tested here, and why.}
The supply-temperature field is the domain of the temperature-propagating formulation, and
two independent obstacles stand between it and a quantitative validation on this network.
The first is formulational: the baseline does not propagate supply temperature, so its
simulated trunk drop is identically zero against a measured $8.2$\,°C, and the far-end and
corridor gates simply do not apply to it. The second is instrumental. Evaluated on the
propagating formulation, the far-end supply-temperature error is $9.21$\,°C, of which
$9.21$\,°C is \emph{bias} -- the mean absolute error and the bias coincide to four decimal
places, so every residual carries the same sign. That is not a model that is wrong but one
that is displaced, by an amount squarely inside the $5$--$15$\,°C offset a three-way mixing
valve produces; the corridor comparison shows the same signature with a bias of
$-7.73$\,°C. Neither obstacle is removable by a better pipe model, and the first is a
deliberate property of the formulation whose cost results we report.

This is the concrete form of an argument that runs through the paper. A model resolved more
finely than its metering cannot have the extra resolution verified, and a model that does
not represent a quantity cannot be tested on it. The conclusions are therefore anchored to
the annual network loss, which both formulations agree on and which the measurements can
resolve, rather than to a point temperature field that neither the instrumentation nor the
baseline can support.

\paragraph{Scope of the temperature evaluation.}
One difference from the original submission is worth stating, since the figures differ. The
original reported a mean supply-temperature error of $1.32$\,°C over six of the fourteen
nodes across a 744-hour winter window, excluding four nodes used in the coefficient
calibration and four with known measurement problems. The figure quoted above is the far-end
node over the full annual record, without selection. Both are correct for what they measure;
the second is the more conservative view and the one we prefer, since the far end is where
the mixing valve runs closest to fully open and the annual record admits no favourable
window. That the earlier selection was not a favourable one is shown in
Appendix~\ref{app:per_node_mae}: the mean across all fourteen nodes, $1.68$\,°C, lies
between the validation-node and calibration-node group means, and two of the excluded nodes
are the best-performing in the network.

\paragraph{Flow, and the per-step balance.}
The source flow shows a mean absolute percentage error of \textbf{33.8}\,\%, and
disaggregating by load band explains it rather than excusing it. The absolute error is
near-constant across bands at roughly $12$--$13$\,m$^3$/h, so the percentage is governed by
the denominator: the same physical error is $46$\,\% of a $30.9$\,m$^3$/h mean flow below a
quarter of peak load and $13$\,\% of a $102.6$\,m$^3$/h mean flow at half to three-quarters
load. Low-load hours are $60$\,\% of the year but carry little energy, which is why a large
relative error there is compatible with an annual energy match of $1.23$\,\%. The top band
is the exception in both directions, with the largest absolute error at $25.5$\,m$^3$/h and
$23$\,\% relative, but comprises 56 hours; we report it rather than absorb it into an
average. The per-step energy balance does not meet its gate, at a mean of $6.4$\,\%: this is
an hour-by-hour closure statistic and is dominated by the same low-flow hours, whereas the
annual closure of $1.23$\,\% -- the quantity the cost conclusions depend on -- is well
inside the gate.

A percentage error can hide either of two quite different failures: a model that mistakes
the level, or one that fails to follow the variation. Comparing first differences separates
them, because a fixed offset cancels. The model tracks the flow level with a correlation of
\textbf{0.91} and its day-to-day change with \textbf{0.80}; for demand the corresponding
figures are \textbf{0.93} and \textbf{0.89}. At hourly resolution the correlation of
differences falls to $0.24$, which is the metering rather than the model: the record is
sampled at fifteen-minute intervals and resampled to hourly, so hour-to-hour differences are
dominated by sampling noise. The aggregate flow error is therefore a level-and-denominator
effect at low load, not a failure to reproduce how the network actually varies -- which is
the property the dispatch results depend on, since cost accumulates over days and seasons
rather than over individual hours.

\paragraph{Post-upgrade assets and structural checks.}
The heat pump, electrode boiler and storage were installed after the measurement record, so
their dispatch is verified against necessary conditions rather than measurements:
coefficient of performance within its plausible band, storage state of charge inside its
operating limits, no simultaneous charge and discharge. Two structural checks close the
section. The nonlinear reference with its quadratic constraints deactivated reproduces the
baseline objective to better than $0.01$\,\%, and the relaxation property
$\mathrm{cost}(\CP) \le \mathrm{cost}(\Lone)$ holds on every one of the 135 synthetic
instances; a violation of either would indicate a formulation error rather than a finding.
